\origpage{256}{ROBERTSON AND SEYMOUR}
\noindent problem until 1979, when Archdeacon, Glover, Huneke, and Wang [1, 3] determined $T(\Sigma)$ when $\Sigma$ is the projective plane. They found that in that case $T(\Sigma)$ has 103 members, up to isomorphism.

One can show that for higher surfaces $\Sigma$, $T(\Sigma)$ has a large number of non-isomorphic members, and it is not clear even that it has only finitely many. The question of the finiteness of $T(\Sigma)$ up to isomorphism was raised by Erd\H{o}s in the 1930's, but again there was little progress until recently. In 1980 Archdeacon and Huneke [2] proved finiteness for all non-orientable surfaces $\Sigma$. The orientable cases have remained open (except for the torus, for which R. Bodendiek and K. Wagner have recently announced a proof of finiteness). In this paper we complete the solution to Erd\H{o}s' problem---we show that for any surface $\Sigma$, $T(\Sigma)$ is indeed finite up to isomorphism.

This is proved as a consequence of a more general ``well-quasi-ordering'' result about graph minors. A graph $H$ is a minor of a graph $G$ if $H$ can be obtained from a subgraph of $G$ by edge-contraction. The ``minor'' relation is different from topological containment---if $H$ is topologically contained in $G$ then it is isomorphic to a minor of $G$, but the converse is not necessarily true, as is easily seen. However, Erd\H{o}s' problem may be reformulated in terms of minors; for, up to isomorphism, $T(\Sigma)$ is finite if and only if the class of minor-minimal graphs which cannot be drawn in $\Sigma$ is finite. (We prove this fact in Section 2.)

We approach the ``minor'' formulation of Erd\H{o}s' problem by means of an idea of Wagner which dates from the 1960's. Let us choose one representative from each isomorphism class of the class of minor-minimal graphs which cannot be drawn in $\Sigma$, and let the class of graphs we obtain be $A$. Now $A$ has the desirable property that it is an antichain, that is, no member of $A$ is isomorphic to a minor of another. (This follows from the fact that the members of $A$ are minor-minimal with a certain property, viz. not being drawable in $\Sigma$.) Wagner made the conjecture that \textit{every} antichain of graphs is finite.

Although we do not prove Wagner's conjecture here, we shall prove the following special case which is strong enough to settle Erd\H{o}s' problem. (We hope to prove Wagner's conjecture in full in a later paper.)

\stmt
(1.1) \textit{For any number $g$, every antichain of graphs all of genus $\le g$ is finite.}

[The genus of a graph is the minimum genus (orientable or non-orientable) of the surfaces in which it can be drawn.]

This settles Erd\H{o}s' problem; for if $\Sigma$ is orientable with genus $g$, it is easy to see that all members of $A$ have genus at most $g+1$, and so by (1.1), $A$ is finite. (The non-orientable case is similar.)

Now for any number $g$ there is a finite list of surfaces such that every
